% Generated draft: requires mathematical and source review.
\documentclass[11pt]{article}
\usepackage[margin=25mm]{geometry}
\usepackage{fontspec}
\setmainfont{Latin Modern Roman}
\usepackage{amsmath,amssymb,mathtools,unicode-math}
\setmathfont{Latin Modern Math}
\usepackage{enumitem,needspace,xurl,hyperref}
\setlength{\parindent}{0pt}
\setlength{\parskip}{5pt}
\setlength{\emergencystretch}{3em}
\newcommand{\N}{\mathbb{N}}
\newcommand{\Z}{\mathbb{Z}}
\newcommand{\Q}{\mathbb{Q}}
\newcommand{\R}{\mathbb{R}}
\newcommand{\C}{\mathbb{C}}
\newcommand{\F}{\mathbb{F}}
\newcommand{\A}{\mathbb{A}}
\newcommand{\PP}{\mathbb{P}}
\newcommand{\E}{\mathbb{E}}
\newcommand{\eps}{\varepsilon}
\newcommand{\dd}{\,\mathrm{d}}
\begin{document}
\section*{UnsolvedMath ID: 30004991}
\textbf{Source problem code:} catalogue entry
\subsection{Definitions and mathematical statement}
The only recoverable mathematical setup is:
\[
G\text{ is a second countable locally compact group.}
\]
Second countability means that the topology of \(G\) has a countable base. Local compactness means that the identity has a compact neighbourhood, and hence every point does by translation.

\textbf{Unresolved mathematical statement.} The supplied extraction contains the heading ``A type I conjecture'', the setup above, and the talk title ``Hyperbolic groups of type I'', attributed to Pierre-Emmanuel Caprace, with joint work by Mehrdad Kalantar and Nicolas Monod. It supplies no conjectural conclusion or quantified assertion beyond the setup. In particular, neither the type-I property nor hyperbolicity can be imposed on \(G\) merely from these headings.

The supplied statement-verification note describes this extraction as a section heading and incomplete setup duplicated by a following complete conjecture. That following conjecture is absent from the supplied record. Consequently, no precise conjecture can be recovered here without additional source text. The record's \texttt{open} label and its assessment of a general locally compact-group conjecture are supplied curation, not a complete mathematical formulation or an independently verified status.
\subsection{Short English statement}
The record introduces a second countable locally compact group but omits the conjecture itself. Its type-I and hyperbolic-group headings do not determine the missing hypotheses or conclusion. This draft preserves the surviving setup and requires the missing source passage before a precise problem can be stated.
\subsection{Sources}
\begin{itemize}
\item Local UnsolvedMath registry, permanent ID \(30004991\), source problem code \texttt{catalogue entry}, title ``Type-I Conjecture Section Heading'', in the Oberwolfach Reports Open Problems collection. The supplied record marks the statement \texttt{unrecoverable} and the problem \texttt{open}; these are registry labels.
\item Supplied report citation: \emph{Groups and Dynamics: Topology, Measure, and Borel Structure (2023)}, Oberwolfach Reports, DOI 10.4171/owr/2022/3, pp.~107--169. \url{https://doi.org/10.4171/owr/2022/3}. The record separately gives proposed year 2022. The extracted talk title is ``Hyperbolic groups of type I'', attributed to Pierre-Emmanuel Caprace, joint work with Mehrdad Kalantar and Nicolas Monod. The supplied verification note says that the official report identifies the extraction as a heading and incomplete setup duplicated by a following complete conjecture; this verification claim has not been independently checked here.
\item Caprace--Kalantar--Monod, \emph{A type I conjecture and boundary representations}. \url{https://arxiv.org/abs/2110.00190}.
\item Glasner, \emph{Boundary representations of locally compact hyperbolic groups}. \url{https://doi.org/10.1112/plms.70100}.
\item Literature assessment dated exactly \texttt{2026-08-22}: the supplied curation describes the extraction as incomplete, reports substantial known hyperbolic and tree cases, and labels the general locally compact-group statement unresolved. This is supplied dated literature triage derived from the OWR curation, not an independently verified assessment. Its time-sensitive status and cited claims require independent verification.
\end{itemize}
\end{document}
